% =============================================================================
% 6. FORMAL PROPERTIES
% Target: 2 pages (~1000 words) — THIS IS THE PAPER'S DIFFERENTIATOR FOR A1
% Source: cap4 §4.7–§4.8 (sec:integration, sec:method-properties)
% Status: SKELETON
% =============================================================================

\section{Formal Properties}
\label{sec:properties}

% ---------- §6.1 Pipeline correctness ----------
\subsection{Pipeline Correctness}
\label{sec:prop-correctness}

\begin{theorem}[Pipeline Correctness]
\label{thm:pipeline-correctness}
If each stage satisfies its contract (Table~\ref{tab:contracts}), then
the composed pipeline $\mathcal{F} = \mathcal{S}_4 \circ \mathcal{S}_3
\circ \mathcal{S}_2 \circ \mathcal{S}_1$ produces a valid empirical
forecast distribution $\hat{\mathcal{D}}$ for any input $\mathcal{R}$
satisfying:
(i)~at least one repository contains events from $\ge 2$ data sources;
(ii)~$|\mathcal{L}| \ge n_{\min}^{\text{cases}}$ (default 50);
(iii)~$|S_T| \ge 1$ and $|S_A| \ge 1$.
\end{theorem}

\begin{proof}
By structural induction on stages: each postcondition implies the
precondition of the next stage. (S1$\to$S2) event-log validity;
(S2$\to$S3) sound model and non-empty state space; (S3$\to$S4)
row-stochastic $P$ and non-empty empirical sojourn distributions; S4
terminates either by absorption certainty
(Property~\ref{prop:absorption-certainty}) or by the $K_{\max}$ safety
bound. \qed
\end{proof}

% ---------- §6.2 Statistical consistency ----------
\subsection{Statistical Consistency}
\label{sec:prop-consistency}

\begin{theorem}[Asymptotic Consistency]
\label{thm:consistency}
As $|\mathcal{L}| \to \infty$ and $M_{\text{sim}} \to \infty$:
\begin{enumerate}
    \item[(i)] $\hat{P} \xrightarrow{a.s.} P^*$;
    \item[(ii)] $\hat{F}_T \xrightarrow{d} F_T^*$;
    \item[(iii)] $\hat{Q}_\alpha \to Q_\alpha^*$ at continuity points.
\end{enumerate}
\end{theorem}

\begin{proof}[Proof sketch]
(i) Property~\ref{prop:mle-consistency}.
(ii) Continuous mapping theorem applied to (i) and the
Glivenko--Cantelli convergence $\hat{F}_s \to F_s^*$.
(iii) Quantile continuity. \qed
\end{proof}

% ---------- §6.3 End-to-end convergence ----------
\subsection{End-to-End Distributional Convergence}
\label{sec:prop-convergence}

\begin{definition}[Forecast Calibration]
\label{def:calibration}
For a finite set of nominal levels
$\mathcal{A} \subset (0,1)$,
\begin{equation}
\label{eq:calibration}
\text{CalErr}(\alpha) = |\hat{P}(T \leq \hat{Q}_\alpha) - \alpha|,
\qquad
\text{ACE} = \tfrac{1}{|\mathcal{A}|}
    \sum_{\alpha \in \mathcal{A}} \text{CalErr}(\alpha).
\end{equation}
\end{definition}

\begin{theorem}[End-to-End Convergence]
\label{thm:pipeline-convergence}
Under (a)~ergodicity of the underlying chain, (b)~uniform sojourn
consistency
$\sup_t |\hat{F}_i^{(n)}(t) - F_i^*(t)| \xrightarrow{a.s.} 0$ for each
$i$, and (c)~$K(n) \to \infty$:
\begin{equation}
\label{eq:pipeline-convergence}
\sup_t |\hat{F}_{n,K}(t) - F^*(t)| \xrightarrow{a.s.} 0,
\qquad
\text{ACE}_{n,K} \to 0.
\end{equation}
\end{theorem}

\begin{proof}[Proof sketch]
SLLN for ergodic Markov chains gives entry-wise
$\hat{P}^{(n)}\to P$~\cite{billingsley1961statistical};
Glivenko--Cantelli per state plus union bound across the finite $S$;
continuous mapping on simulated trajectories plus a second
Glivenko--Cantelli application across $K$ replications gives the
uniform convergence. ACE convergence follows.
\end{proof}

% ---------- §6.4 Information-theoretic advantage ----------
\subsection{Information-Theoretic Advantage of Process Awareness}
\label{sec:prop-information}

\begin{theorem}[Information Gain from Process Structure]
\label{thm:information}
\begin{equation}
\label{eq:information-gain}
I(T; \psi) = H(T) - H(T \mid \psi) \ge 0,
\end{equation}
with equality iff $T$ is independent of the trajectory $\psi$.
\end{theorem}

\paragraph{Practical interpretation.}
In real software processes $I(T;\psi) > 0$ because rework-laden paths
have systematically longer lead times: the trajectory $\psi$ carries
information about state-specific delays and revisits that the
marginal $T$ does not. Process-aware Monte Carlo (Stage~4) samples
the joint $(T, \psi)$ and captures this variance decomposition;
throughput-based Monte Carlo collapses $\psi$ into a daily-throughput
scalar and forfeits the information gain.

% ---------- §6.5 Complexity and assumptions ----------
\subsection{Complexity and Assumptions}
\label{sec:prop-complexity}

The per-stage time complexity is summarized in
Table~\ref{tab:complexity}; the total cost is dominated by Stage~4:
\begin{equation}
\label{eq:total-cost}
C_{\text{total}} = \mathcal{O}(|\mathcal{L}|)
                 + \mathcal{O}(M_{\text{sim}} \cdot \bar{K}),
\end{equation}
where $\bar{K}$ is the expected trajectory length, bounded by
$K_{\max}$. Monte Carlo is embarrassingly parallel and admits GPU
offload~\cite{lecuyer2017random}.

\begin{table}[ht]
\centering
\caption{Per-stage time complexity.}
\label{tab:complexity}
\small
\begin{tabular}{@{}lll@{}}
\toprule
\textbf{Stage} & \textbf{Time complexity} & \textbf{Dominant factor} \\
\midrule
$\mathcal{S}_1$: Event Log     & $O(|E^{\text{raw}}| \cdot |\mathcal{K}|)$ & raw events $\times$ strategies \\
$\mathcal{S}_2$: Discovery + Conf. & $O(|\mathcal{L}| \cdot |\mathcal{A}|^2 + |\mathcal{C}| \cdot |M|)$ & alignment cost \\
$\mathcal{S}_3$: Markov         & $O(|\mathcal{L}| + |S_T|^3)$ & counting + matrix inversion \\
$\mathcal{S}_4$: Monte Carlo    & $O(M_{\text{sim}} \cdot \bar{K})$ & simulations $\times$ path length \\
$\mathcal{S}_{\text{ML}}$: Training & $O(|\mathcal{C}| \cdot p \cdot T)$ & cases $\times$ features $\times$ trees \\
\bottomrule
\end{tabular}
\end{table}

The theoretical results rest on five explicit assumptions, listed in
Table~\ref{tab:assumptions} together with the consequence of
violation and the mitigation adopted by PATHCAST.

\begin{table}[ht]
\centering
\caption{Assumptions, violation consequences and mitigations.}
\label{tab:assumptions}
\small
\begin{tabular}{@{}p{2.6cm} p{4.6cm} p{4.6cm}@{}}
\toprule
\textbf{Assumption} & \textbf{Violation consequence} & \textbf{Mitigation} \\
\midrule
A1 First-order Markov            & Underestimates rework-conditioned delays & State augmentation; ML features \\
A2 Stationarity                   & Stale forecasts after process change      & Sliding window; drift detection \\
A3 Case independence              & Ignores queueing/resource contention      & Acknowledged limitation \\
A4 Sojourn-history independence   & Underestimates reworked-item delays       & State augmentation; ML features \\
A5 Complete case observation      & Bias if long cases excluded               & Include right-censored cases \\
\bottomrule
\end{tabular}
\end{table}
